\documentclass{article}
\usepackage{amsmath}
\usepackage{amssymb}

% Define the viz environment (consumes optional args so they don't render in PDF)
\newenvironment{viz}[1][]{\begin{equation}}{\end{equation}}

\begin{document}

\title{LaTeX Visualiser Demo}
\author{Demo}
\date{}


The two-sided cone $S=\{(x,y,z)\in\mathbb{R}^3:x^2+y^2=z^2,-2\leq z\leq2\}$ is parametrized by the cylindrical coordinates $\max G:[0,2\pi]\times[-2,2]\to\mathbb{R}^3$ given by
$$G(s,t)=(t\cos s,t\sin s,t).$$

$\textbf{Notice that}$

$$\begin{aligned}&\partial_1G(s,t)=(-t\sin s,t\cos s,0),\\&\partial_2G(s,t)=(\cos s,\sin s,1).\end{aligned}$$
You can verify that the map $G$ is smooth, and $\{\partial_1G(s,t),\partial_2G(s,t)\}$ is linearly dependent if and only if $s\in(0,2\pi)$ and $t=0.$ Play with this Math3D demo to



view this phenomenon.
Thus, $G$ is not regular so you cannot conclude whether $S$ is a surface. In fact, $S$ is not a surface but it is not easy to formally prove as you need to exclude all possible parametrizations You can, however, write $S$ as a union of two surfaces.
$\bullet$ Define the upper cone

$$S_1=\{(x,y,z)\in\mathbb{R}^3:x^2+y^2=z^2,0\leq z\leq2\}.$$

The map $G_1:U_1\to\mathbb{R}^3$ given by $U_1=[0,2\pi]\times[0,2]$ and $G_1(s,t)=(t\cos s,t\sin s,t)$

is a simple smooth regular parametrization of $S_1$ and hence $S_1$ is a surface.

\begin{viz}[type=curve, label=AI Suggested Curve, xrange=-6:6, yrange=-2:2]
y = sin(x)
\end{viz}

$\bullet$ Define the lower cone

$$S_2=\{(x,y,z)\in\mathbb{R}^3:x^2+y^2=z^2,-2\leq z\leq0\}.$$

The map $G_2:U_2\to\mathbb{R}^3$ given by $U_2=[0,2\pi]\times[-2,0]$ and $G_2(s,t)=(t\cos s,t\sin s,t)$
is a simple smooth regular parametrization of $S_2$ and hence $S_2$ is a surface.
Thus, $S=S_1\cup S_2$ and, more importantly, the intersection $S_1\cap S_2$ is equal to the image of the
$\textbf{domain boundaries of }G_1$ and $G_2$, namely

$$S_1\cap S_2=G_1(\partial U_1)\cap G_2(\partial U_2).$$

The surfaces $S_1$ and $S_2$ are "glued along their parametrization's domain boundaries"

% courtesy of MAT237 text book

\end{document}
